\subsection{Аналіз класичної інтерполяції Лагранжа}

Вихідним пунктом для побудови апарату масштабування зображень є класична задача поліноміальної інтерполяції. Найбільш відомим підходом до її розв’язання є використання інтерполяційного полінома у формі Лагранжа, який для набору вузлів $\{x_j\}_{j=0}^n$ та відповідних значень функції $\{f_j\}_{j=0}^n$ записується як:
\begin{equation}\tag{1.1}
	p(x) = \sum_{j=0}^{n} f_j \ell_j(x),
\end{equation}
де $\ell_j(x)$ — базисні поліноми Лагранжа (фундаментальні поліноми), що мають вигляд:
\begin{equation}\tag{1.2}
	\ell_j(x) = \prod_{k=0, k \neq j}^{n} \frac{x - x_k}{x_j - x_k}.
\end{equation}

З теоретичної точки зору, поліном Лагранжа є єдиним поліномом ступеня $n$, який точно збігається з функцією у заданих вузлах. Проте при програмній реалізації цього підходу для задач обробки зображень виникає ряд суттєвих обмежень:

\begin{itemize}
	\item \textbf{Висока обчислювальна складність:} Пряме обчислення виразу (1.2) для кожного нового значення $x$ вимагає $O(n)$ операцій, що при підсумовуванні в (1.1) призводить до загальної складності $O(n^2)$. Для зображень високої роздільної здатності це робить процес масштабування неприпустимо повільним.
	\item \textbf{Числова нестабільність:} Використання базисних поліномів (1.2) у стандартній формі є вкрай чутливим до похибок округлення при великих значеннях $n$, що часто призводить до втрати точності.
	\item \textbf{Відсутність гнучкості:} Додавання нових вузлів у сітку вимагає повного перерахунку всього полінома, що унеможливлює динамічну зміну масштабу.
\end{itemize}

Саме ці недоліки класичної форми Лагранжа обґрунтовують необхідність переходу до барицентричної інтерполяційної форми. Вона дозволяє значно спростити обчислювальний процес та уникнути проблем із втратою точності, властивих класичним поліномам Лагранжа, забезпечуючи при цьому високу ефективність наближення.